% ============================================================
% DRAFT NOTE (forecaster)
% Contains: What the forecaster is; the raw material and the causal contract
% (protocol section 5); feature engineering and ablation; model class and
% frozen selection; calibration; inference cost. Comparators live in Setup;
% interpretation of what the model looks at lives in Results.
% Written now: Paragraphs 1, 2, and the structure of 4-6 (protocol-locked).
% Pending (agent decisions): engineered features and ablation outcome;
% whether the frozen divergence outputs were used; selected model class and
% hyperparameters; calibration method; measured per-token cost.
% Sources: docs/implementation/quality_risk_protocol.md sections 5 and 7;
% current-state-damage-v1 protocol_lock.json (feature names, nested scheme)
% ============================================================
\section{The Forecaster}
\label{sec:forecaster}

HERALD is a classifier. At each decoding step \step{} of a compressed run it
maps the causally available information \info{} to an estimate of
$\risk(\step)$, the probability that the run ends damaged. It is trained on
every compressed token row of the development prompts, each row carrying its
run's label \dmg, so the same model serves every position and every
compression action; position and action are inputs, not separate models.
Because the label is shared by all rows of a run, early rows are noisy
training examples by construction: the evidence at \step{} may not yet
determine the outcome. That is the intended difficulty. The forecaster
learns how much the evidence available at each point actually says about the
outcome, and the earliness curve of Section~\ref{sec:results} reports the
answer.

\paragraph{Raw material and the causal contract.}
Every input is derived from three sources. The compression action
$\action = (\press, \ratio)$, encoded as compressor identity, ratio, and their
combination, which carries the base rate of damage for that configuration.
The absolute position \step. And the twelve per-token statistics of the
compressed model's own next-token distribution described in
Section~\ref{sec:corpus}: entropy, top-1 and top-5 probability, the number of
plausible alternatives, average log-probability, entropy change and KL
divergence from the previous step, top-10 overlap with the previous step,
effective vocabulary size, tail probability mass, and logit range. A feature
is admissible if it can be recomputed from rows of the same run with
position at most \step; nothing else enters. This excludes the teacher-forced
divergence oracles, because they require the full-cache model; the task
scores and catastrophe flags, because they are the label or its
correlates; relative progress and the final length, because they encode how
the run ends; and the task identity, which is used to stratify results but
never to predict them. The contract is what makes the forecaster zero-cost:
at inference time it adds a fixed-size update of running statistics over
logits the decoder has already produced, and one model evaluation.

\paragraph{Feature engineering.}
On top of the raw material we build causal transformations: trailing means
and standard deviations over 8- and 32-token windows and exponential moving
averages at matching half-lives for the principal signals, following the
release, and \todo{describe the engineered features that survived
development: differences, ratios, interactions, cumulative statistics, shape
features of the recent trajectory, as reported by the agent}. \todo{State
whether the four outputs of the frozen divergence forecaster of
current-state-damage-v1 were included, and whether they survived.} Feature
selection used development folds only, and the final feature set is frozen
by code hash before confirmation. Table~\ref{tab:ablation} groups the
features as instantaneous state, step-to-step dynamics, windowed history,
and action-and-position, and reports development performance with each group
removed; the interpretation of what the forecaster is looking at is deferred
to Section~\ref{sec:results}.

\begin{table}[t]
\centering
\small
\caption{Feature-group ablation on development folds (prompt-grouped
out-of-fold). \todo{Fill from the agent's development report; choose the
checkpoint and metric consistent with Section~\ref{sec:setup}.}}
\label{tab:ablation}
\begin{tabular}{lcc}
\toprule
Features removed & AUROC at $\step{=}32$ & Log loss at $\step{=}32$ \\
\midrule
None (full model)            & \todo{v} & \todo{v} \\
Instantaneous state          & \todo{v} & \todo{v} \\
Step-to-step dynamics        & \todo{v} & \todo{v} \\
Windowed history             & \todo{v} & \todo{v} \\
All sensors (action-clock)   & \todo{v} & \todo{v} \\
\bottomrule
\end{tabular}
\end{table}

\paragraph{Model class and selection.}
The first candidate is a gradient-boosted tree classifier over the frozen
feature set, with \step{} as an ordinary input. The challenger is a causal
sequence model over the raw per-token signals, a temporal convolutional
network with a sigmoid head, which can in principle learn its own history
features. Both are trained under the same nested prompt-grouped scheme:
within the development prompts, one fold provides early stopping, three
train, and the outer fold is held out, so hyperparameters are never chosen on
data they are scored on. The selection rule was frozen before confirmation:
if both candidates clear the development gates, the tree model is selected
unless the sequence model's macro loss is better with a one-sided 95\%
prompt-clustered bound. \todo{Report the selected model, its capacity and
hyperparameters, and the development margin between the two.}

\paragraph{Calibration.}
A risk estimate is only useful to a decision if its probabilities mean what
they say. The selected model's scores are calibrated with
\todo{isotonic regression or Platt scaling; state which} fit on
out-of-fold development scores only, and the calibrated map is frozen with
the model. Calibration quality is a gate, not a diagnostic
(Section~\ref{sec:setup}).

\paragraph{Cost at inference.}
Per token, the forecaster updates a fixed number of running statistics and
evaluates one model; the cost is independent of sequence length and of the
compressor. \todo{Report measured per-token latency of the feature update
and model evaluation on CPU and, if measured, alongside a decode step on
GPU.}
